# Введение

Составление расписания занятий — одна из тех задач университетской жизни, которая кажется простой ровно до тех пор, пока не приходится решать её на практике. На первый взгляд нужно всего лишь расставить пары так, чтобы преподаватель не оказался одновременно в двух аудиториях, а группа не получила две дисциплины в одно и то же время. Но как только в задачу добавляются реальные ограничения — вместимость аудиторий, специализация кабинетов под лабораторные работы, нагрузка преподавателей, желание не растягивать день большими окнами между парами — простая на словах задача превращается в комбинаторную головоломку, ручное решение которой занимает у методистов вуза дни, а иногда недели в начале каждого семестра.

Это не локальная проблема одного учебного заведения. Задача составления расписания относится к классу NP-трудных задач: с ростом числа дисциплин, групп, преподавателей и аудиторий количество возможных вариантов расписания растёт настолько быстро, что полный перебор всех вариантов перестаёт быть осуществимым на практике уже при умеренных объёмах данных. Это делает задачу интересной не только с практической, но и с алгоритмической точки зрения: она оказывается удобным полигоном для сравнения разных подходов к решению NP-трудных задач — от точных методов до эвристик, не гарантирующих оптимальность, но способных найти хорошее решение за разумное время.

Актуальность темы для практики очевидна: автоматизация составления расписания напрямую экономит время методистов и снижает число ошибок, неизбежных при ручной расстановке сотен пар. Актуальность для учебной и исследовательской работы — в том, что задача составления расписания одновременно простая для понимания (любой человек интуитивно понимает условия и ограничения) и достаточно сложная, чтобы на ней можно было содержательно сравнить разные алгоритмические парадигмы: эволюционные методы, методы локального поиска, точное программирование ограничений. Это даёт возможность не просто реализовать один алгоритм, а провести осмысленное сравнение нескольких принципиально разных подходов на одной и той же постановке задачи.

Целью данной работы является разработка и сравнительный анализ нескольких алгоритмов автоматического построения расписания учебных занятий, способных учитывать как обязательные (жёсткие), так и желательные (мягкие) ограничения, присущие реальному учебному процессу.

Для достижения этой цели были поставлены следующие задачи:

Во-первых, изучить теоретические основы задачи составления расписания и существующие подходы к её решению, чтобы выбрать алгоритмы, наиболее подходящие для рассматриваемой постановки задачи. Во-вторых, спроектировать единую модель данных и единый критерий оценки качества расписания, на основе которых можно реализовать и затем честно сравнить между собой разные алгоритмы. В-третьих, реализовать четыре алгоритма построения расписания — генетический алгоритм, имитацию отжига, поиск с запретами и программирование ограничений (CP-SAT) — на основе этой единой модели. В-четвёртых, провести экспериментальное сравнение реализованных алгоритмов по нескольким критериям: качество найденного решения, скорость работы, устойчивость к разной структуре исходных данных. В-пятых, сформулировать практические выводы о применимости каждого из подходов в зависимости от размера и характера задачи.

Методологическую базу работы составили анализ научной и технической литературы по теории расписаний и методам комбинаторной оптимизации, проектирование и реализация программного обеспечения на языке Python с использованием библиотек DEAP (генетические алгоритмы) и OR-Tools (CP-SAT), а также вычислительный эксперимент: серия прогонов всех четырёх алгоритмов на синтетических наборах данных различного размера и сложности с последующим статистическим сравнением результатов.

Работа состоит из трёх глав. В первой главе рассматриваются теоретические основы задачи составления расписания: формальная постановка задачи, понятие жёстких и мягких ограничений, а также краткий обзор существующих программных решений в этой области. Во второй главе подробно описывается каждый из четырёх реализованных алгоритмов — принципы его работы, особенности конкретной реализации, проблемы, возникшие в процессе разработки, и выводы по применимости каждого алгоритма, — а также приводится сравнительный анализ всех четырёх подходов по результатам вычислительного эксперимента. В заключении подводятся итоги проделанной работы и обозначаются возможные направления её дальнейшего развития.
